\chapter{2013年北京卷（理）}

\section{单选题}

\begin{problem}
已知集合 \(A = \{-1, 0, 1\}\)，\(B = \{x \mid -1 \leqslant x < 1\}\)，则 \(A \cap B =\)（\quad）

\choices
    {\(\{0\}\)}
    {\( \{-1,0\} \)}
    {\(\{0,1\}\)}
    {\( \{-1, 0, 1\} \)}
\end{problem}

\begin{answer}
B
\end{answer}

\begin{solution}
集合 \(B\) 表示区间 \([-1,1)\)，其中包含 \(-1\) 而不包含 \(1\). 因此 \(A\) 中同时属于 \(B\) 的元素为 \(-1\) 和 \(0\)，即 \(A\cap B=\{-1,0\}\). 故选 B.
\end{solution}

\begin{problem}
在复平面内，复数 \((2 - \i)^2\) 对应的点位于（\quad）

\choices
    {第一象限}
    {第二象限}
    {第三象限}
    {第四象限}
\end{problem}

\begin{answer}
D
\end{answer}

\begin{solution}

\((2-\i)^2=4-4\i+\i^2=3-4\i\)，其实部为正、虚部为负，所以对应点位于第四象限. 故选 D.
\end{solution}

\begin{problem}
“\(\varphi = \pi\)”是“曲线 \(y = \sin (2x + \varphi)\) 过坐标原点”的（\quad）

\choices
    {充分而不必要条件}
    {必要而不充分条件}
    {充分必要条件}
    {既不充分也不必要条件}
\end{problem}

\begin{answer}
A
\end{answer}

\begin{solution}

曲线过坐标原点的条件是令 \(x=0\) 时 \(y=0\)，即 \(\sin\varphi=0\)，所以 \(\varphi=k\pi\)，其中 \(k\in\Z\). 当 \(\varphi=\pi\) 时确有 \(\sin\pi=0\)，故该条件是充分条件；但例如 \(\varphi=0\) 时曲线也过原点，故它不是必要条件. 故选 A.
\end{solution}

\begin{problem}
执行如图所示的程序框图，输出的 \(S\) 值为（\quad）

\bitmapfigure[max width=0.32\linewidth]{img/2013/beijing_science/q4_fig1.png}

\choices
    {1}
    {\(\frac{2}{3}\)}
    {\(\frac{13}{21}\)}
    {\(\frac{610}{987}\)}
\end{problem}

\begin{answer}
C
\end{answer}

\begin{solution}

初始时 \(i=0\)，\(S=1\). 第一次循环后
\(S=\dfrac{1^2+1}{2\cdot1+1}=\dfrac23\)，且 \(i=1\)；第二次循环后
\(S=\dfrac{(\frac23)^2+1}{2\cdot\frac23+1}=\dfrac{\frac49+1}{\frac73}=\dfrac{13}{21}\)，且 \(i=2\). 此时满足 \(i\geqslant2\)，程序停止，故选 C.
\end{solution}

\begin{problem}
函数 \( f(x) \) 的图象向右平移 1 个单位长度，所得图象与曲线 \( y = \e^{x} \) 关于 \( y \) 轴对称，则 \( f(x) = \)（\quad）

\choices
    {\(\e^{x + 1}\)}
    {\(\e^{x - 1}\)}
    {\(\e^{-x + 1}\)}
    {\(\e^{-x - 1}\)}
\end{problem}

\begin{answer}
D
\end{answer}

\begin{solution}

将 \(y=\e^x\) 关于 \(y\) 轴对称，得到曲线 \(y=\e^{-x}\). 函数 \(f(x)\) 的图象向右平移 \(1\) 个单位后为 \(y=f(x-1)\)，所以 \(f(x-1)=\e^{-x}\). 令 \(u=x-1\)，则 \(f(u)=\e^{-(u+1)}\)，从而 \(f(x)=\e^{-x-1}\). 故选 D.
\end{solution}

\begin{problem}
若双曲线 \(\frac{x^2}{a^2} - \frac{y^2}{b^2} = 1\) 的离心率为 \(\sqrt{3}\)，则其渐近线方程为（\quad）

\choices
    {\(y = \pm 2x\)}
    {\(y = \pm \sqrt{2} x\)}
    {\(y = \pm \frac{1}{2} x\)}
    {\(y = \pm \frac{\sqrt{2}}{2} x\)}
\end{problem}

\begin{answer}
B
\end{answer}

\begin{solution}

双曲线满足 \(c^2=a^2+b^2\)，由离心率 \(\dfrac ca=\sqrt3\) 得 \(c^2=3a^2\)，所以 \(b^2=2a^2\)，即 \(\dfrac ba=\sqrt2\). 双曲线的渐近线为 \(y=\pm\dfrac ba x\)，故方程为 \(y=\pm\sqrt2 x\). 故选 B.
\end{solution}

\begin{problem}
直线 \(l\) 过抛物线 \(C: x^{2} = 4y\) 的焦点且与 \(y\) 轴垂直，则 \(l\) 与 \(C\) 所围成的图形的面积等于（\quad）

\choices
    {\(\frac{4}{3}\)}
    {2}
    {\( \frac{8}{3} \)}
    {\(\frac{16\sqrt{2}}{3}\)}
\end{problem}

\begin{answer}
C
\end{answer}

\begin{solution}

抛物线 \(x^2=4y\) 的焦点为 \((0,1)\)，所以直线 \(l\) 的方程为 \(y=1\). 联立 \(y=1\) 与 \(y=\dfrac{x^2}{4}\)，得交点的横坐标为 \(x=\pm2\). 所围图形的面积为
\(S=\int_{-2}^{2}\left(1-\dfrac{x^2}{4}\right)\,\mathrm{d}x=\dfrac83\). 故选 C.
\end{solution}

\begin{problem}
设关于 \(x\)，\(y\) 的不等式组 \(\begin{cases}
2x - y + 1 > 0,\\
x + m < 0,\\
y - m > 0
\end{cases}\) 表示的平面区域内存在点 \(P(x_0, y_0)\)，满足 \(x_0 - 2y_0 = 2\)，则 \(m\) 的取值范围是（\quad）

\choices
    {\(\left(-\infty, \frac{4}{3}\right)\)}
    {\(\left(-\infty, \frac{1}{3}\right)\)}
    {\(\left(-\infty, -\frac{2}{3}\right)\)}
    {\(\left(-\infty, -\frac{5}{3}\right)\)}
\end{problem}

\begin{answer}
C
\end{answer}

\begin{solution}

由 \(x_0-2y_0=2\) 得 \(y_0=\dfrac{x_0}{2}-1\). 将其代入不等式组的三个条件，分别得到
\(x_0>-\dfrac43\)、\(x_0>2m+2\) 和 \(x_0<-m\). 因此存在这样的 \(x_0\) 的充要条件是
\(\max\left\{-\dfrac43,2m+2\right\}<-m\)，即同时满足 \(m<\dfrac43\) 和 \(m<-\dfrac23\). 故 \(m\in\left(-\infty,-\dfrac23\right)\)，选 C.
\end{solution}

\section{填空题}

\begin{problem}
在极坐标系中，点 \(\left(2, \frac{\pi}{6}\right)\) 到直线 \(\rho \sin \theta = 2\) 的距离等于\fillinblank{}.
\end{problem}

\begin{answer}
\(1\)
\end{answer}

\begin{solution}

在极坐标中，点 \(\left(2,\frac{\pi}{6}\right)\) 的直角坐标为
\((x,y)=\left(2\cos\frac{\pi}{6},2\sin\frac{\pi}{6}\right)=(\sqrt3,1)\). 直线 \(\rho\sin\theta=2\) 即 \(y=2\)，所以点到该直线的距离为 \(\lvert1-2\rvert=1\).
\end{solution}

\begin{problem}
若等比数列 \( \{a_{n}\} \) 满足 \(a_{2} + a_{4} = 20\)，\(a_{3} + a_{5} = 40\)，则公比 \(q =\)\fillinblank{}，前 \(n\) 项和 \( S_{n} = \)\fillinblank{}.
\end{problem}

\begin{answer}
\(2\)，\(2^{n+1} - 2\)
\end{answer}

\begin{solution}

由 \(a_3=qa_2\)，\(a_5=qa_4\)，有
\(a_3+a_5=q(a_2+a_4)\)，所以 \(40=20q\)，得 \(q=2\). 再由 \(a_2+a_4=a_2(1+q^2)=20\)，得 \(5a_2=20\)，即 \(a_2=4\)，从而 \(a_1=2\). 因此
\(S_n=a_1\dfrac{q^n-1}{q-1}=2(2^n-1)=2^{n+1}-2\).
\end{solution}

\begin{problem}
如图，\(A\) 为圆 \(O\) 的直径，\(PA\) 为圆 \(O\) 的切线，\(PB\) 与圆 \(O\) 相交于 \(D\)，若 \(PA = 3\)，\(PD:DB = 9:16\)，则 \(PD =\)\fillinblank{}，\(AB =\)\fillinblank{}.
\bitmapfigure[max width=0.45\linewidth]{img/2013/beijing_science/q11_fig1.png}
\end{problem}

\begin{answer}
\(\frac{9}{5}\)，\(4\)
\end{answer}

\begin{solution}

由切割线定理，\(PA^2=PD\cdot PB\). 设 \(PD=9t\)，\(DB=16t\)，则 \(PB=25t\)，于是
\(9=9t\cdot25t\)，解得 \(t=\dfrac15\)，所以 \(PD=\dfrac95\)，\(PB=5\). 又 \(OA\perp PA\)，且 \(O\)，\(A\)，\(B\) 共线，所以 \(\triangle PAB\) 是直角三角形，从而
\(AB=\sqrt{PB^2-PA^2}=\sqrt{5^2-3^2}=4\).
\end{solution}

\begin{problem}
将序号分别为 1， 2， 3， 4， 5 的 5 张参观券全部分给 4 人，每人至少 1 张，如果分给同一人的 2 张参观券连号，那么不同的分法种数是\fillinblank{}.
\end{problem}

\begin{answer}
\(96\)
\end{answer}

\begin{solution}

由于 \(5\) 张券分给 \(4\) 人且每人至少 \(1\) 张，人数分配必为一人得 \(2\) 张、其余三人各得 \(1\) 张. 先选得两张券的人，有 \(4\) 种；这两张券必须是相邻编号，共有 \((1,2)\)、\((2,3)\)、\((3,4)\)、\((4,5)\) 四种；剩下三张券分给其余三人，有 \(3!\) 种. 因此不同分法数为 \(4\cdot4\cdot3!=96\).
\end{solution}

\begin{problem}
向量 \(\bs{a}\)，\(\bs{b}\)，\(\bs{c}\) 在正方形网格中的位置如图所示，若 \(\bs{c} = \lambda \bs{a} + \mu \bs{b}\)（\(\lambda\)，\(\mu \in \R\)），则 \(\frac{\lambda}{\mu} =\)\fillinblank{}.
\bitmapfigure[max width=0.40\linewidth]{img/2013/beijing_science/q13_fig1.png}
\end{problem}

\begin{answer}
\(4\)
\end{answer}

\begin{solution}

以网格的水平、竖直方向为坐标轴，由图中各向量的起点和终点可取
\(\bs{a}=(-1,1)\)，\(\bs{b}=(6,2)\)，\(\bs{c}=(-1,-3)\). 由
\(\bs{c}=\lambda\bs{a}+\mu\bs{b}\)，比较两个坐标得
\(-\lambda+6\mu=-1\)，\(\lambda+2\mu=-3\). 两式相加得 \(\mu=-\dfrac12\)，代回得 \(\lambda=-2\)，所以
\(\dfrac{\lambda}{\mu}=4\).
\end{solution}

\begin{problem}
如图，在棱长为 2 的正方体 \(ABCD - A_{1}B_{1}C_{1}D_{1}\) 中，\(E\) 为 \(BC\) 的中点，点 \(P\) 在线段 \(D_{1}E\) 上，点 \(P\) 到直线 \(CC_{1}\) 的距离的最小值为\fillinblank{}.
\bitmapfigure[max width=0.52\linewidth]{img/2013/beijing_science/q14_fig1.png}
\end{problem}

\begin{answer}
\(\frac{2\sqrt{5}}{5}\)
\end{answer}

\begin{solution}

建立空间直角坐标系，令
\(A=(0,0,0)\)，\(B=(2,0,0)\)，\(C=(2,2,0)\)，\(D=(0,2,0)\)，并令竖直方向为 \(z\) 轴，则 \(D_1=(0,2,2)\)，\(C_1=(2,2,2)\). 因为 \(E\) 是 \(BC\) 的中点，所以 \(E=(2,1,0)\). 设
\(P=D_1+t(E-D_1)=(2t,2-t,2-2t)\)，其中 \(0\leqslant t\leqslant1\).

直线 \(CC_1\) 的方程为 \(x=2\)，\(y=2\)，故
\(d^2(P,CC_1)=(2t-2)^2+(2-t-2)^2=5t^2-8t+4=5\left(t-\dfrac45\right)^2+\dfrac45\). 当 \(t=\dfrac45\) 时取得最小值，因此最小距离为
\(\sqrt{\dfrac45}=\dfrac{2\sqrt5}{5}\).
\end{solution}

\section{解答题}

\begin{problem}
在 \(\triangle ABC\) 中，\(a = 3\)，\(b = 2\sqrt{6}\)，\(\angle B = 2\angle A\).
\begin{enumerate}
\item 求 \( \cos A \) 的值.
\item 求 \(c\) 的值.
\end{enumerate}
\end{problem}

\begin{solution}
\begin{enumerate}
\item 由三角形内角和及 \(B=2A\)，有 \(3A<\pi\)，故 \(0<A<\dfrac{\pi}{3}\)，从而 \(\sin A>0\) 且 \(\cos A>0\). 由正弦定理，
\(\dfrac{a}{\sin A}=\dfrac{b}{\sin B}\)，所以
\(\dfrac{3}{\sin A}=\dfrac{2\sqrt6}{\sin2A}=\dfrac{2\sqrt6}{2\sin A\cos A}\). 约去正数 \(\sin A\)，得 \(3\cos A=\sqrt6\)，即 \(\cos A=\dfrac{\sqrt6}{3}\).

\item 由余弦定理，
\(a^2=b^2+c^2-2bc\cos A\)，所以
\(9=24+c^2-2\cdot2\sqrt6\cdot c\cdot\dfrac{\sqrt6}{3}=c^2-8c+24\). 因此
\(c^2-8c+15=0\)，解得 \(c=3\) 或 \(c=5\).

若 \(c=3\)，则 \(a=c\)，所以 \(A=C\). 结合 \(B=2A\) 和内角和，得 \(A=C=\dfrac{\pi}{4}\)，\(B=\dfrac{\pi}{2}\). 于是 \(b^2=a^2+c^2=18\)，但题设 \(b^2=24\)，矛盾，故 \(c\ne3\).

当 \(c=5\) 时，三边 \(3\)，\(2\sqrt6\)，\(5\) 满足三角形不等式，且
\(\cos B=\dfrac{a^2+c^2-b^2}{2ac}=\dfrac{9+25-24}{30}=\dfrac13\). 另一方面，
\(\cos2A=2\cos^2A-1=2\cdot\dfrac{6}{9}-1=\dfrac13\). 又 \(B\)，\(2A\in(0,\pi)\)，余弦函数在该区间上单调递减，所以 \(B=2A\)，该候选值满足题设. 因此 \(c=5\).
\end{enumerate}
\end{solution}

\begin{problem}
下图是某市 3 月 1 日至 14 日的空气质量指数趋势图，空气质量指数小于 100 表示空气质量优良，空气质量指数大于 200 表示空气重度污染. 某人随机选择 3 月 1 日至 3 月 13 日中的某一天到达该市，并停留 2 天.

\begin{enumerate}
\item 求此人到达当日空气重度污染的概率；
\item 设 \(X\) 是此人停留期间空气质量优良的天数，求 \(X\) 的分布列与数学期望；
\item 由图判断从哪天开始连续三天的空气质量指数方差最大（结论不要求证明）
\end{enumerate}

\FigureLayoutDeclare{img/2013/beijing_science/q16_fig1.png}{16.0cm}{61bp 40bp 0bp 21bp}
\bitmapfigure[max width=0.78\linewidth]{img/2013/beijing_science/q16_fig1.png}
\end{problem}

\begin{solution}
\begin{enumerate}
\item 到达日从 3 月 1 日至 3 月 13 日共有 \(13\) 个等可能选择. 图中空气质量指数大于 \(200\) 的日期为 3 月 5 日和 3 月 8 日，因此到达当日空气重度污染的概率为 \(\dfrac{2}{13}\).

\item 空气质量指数小于 \(100\) 的日期为 3 月 1 日、2 日、3 日、7 日、12 日、13 日和 14 日. 按到达日从 3 月 1 日至 3 月 13 日排列，停留两天中空气质量优良的天数依次为
\(2\)，\(2\)，\(1\)，\(0\)，\(0\)，\(1\)，\(1\)，\(0\)，\(0\)，\(0\)，\(1\)，\(2\)，\(2\). 因此
\begin{center}
\begin{tabular}{c|ccc}
\(X\) & \(0\) & \(1\) & \(2\) \\
\hline
\(P\) & \(\dfrac{5}{13}\) & \(\dfrac{4}{13}\) & \(\dfrac{4}{13}\)
\end{tabular}
\end{center}
\(P(X=0)=\dfrac5{13}\)，\(P(X=1)=\dfrac4{13}\)，\(P(X=2)=\dfrac4{13}\). 数学期望为
\(E(X)=0\cdot\dfrac5{13}+1\cdot\dfrac4{13}+2\cdot\dfrac4{13}=\dfrac{12}{13}\).

\item 对连续三天的空气质量指数 \(x\)，\(y\)，\(z\)，其方差为
\(\sigma^2=\dfrac{(x-y)^2+(y-z)^2+(z-x)^2}{9}\). 由图读出 3 月 1 日至 14 日的指数依次为
\(86\)，\(25\)，\(57\)，\(143\)，\(220\)，\(160\)，\(40\)，\(217\)，\(160\)，\(121\)，\(158\)，\(86\)，\(79\)，\(37\). 从 3 月 1 日开始的 12 个三日组对应的 \(9\sigma^2\) 依次为
\(5586\)，\(22344\)，\(39894\)，\(9818\)，\(50400\)，\(48978\)，\(48978\)，\(13986\)，\(2894\)，\(7778\)，\(11474\)，\(4214\). 最大值为 \(50400\)，对应 3 月 5 日、6 日、7 日，所以从 3 月 5 日开始连续三天的空气质量指数方差最大.
\end{enumerate}
\end{solution}

\begin{problem}
如图，在三棱柱 \(ABC - A_{1}B_{1}C_{1}\) 中，\(AA_{1}C_{1}C\) 是边长为 4 的正方形. 平面 \(ABC \perp\) 平面 \(AA_{1}C_{1}C\)，\(AB = 3\)，\(BC = 5\).
\begin{enumerate}
\item 求证： \(AA_{1} \perp\) 平面 \(ABC\)；
\item 求二面角 \(A_{1} - BC_{1} - B_{1}\) 的余弦值；
\item 证明：在线段 \(BC_{1}\) 上存在点 \(D\)，使得 \(AD \perp A_{1}B\)，并求 \(\frac{BD}{BC_{1}}\) 的值.
\end{enumerate}
\bitmapfigure[max width=0.56\linewidth]{img/2013/beijing_science/q17_fig1.png}
\end{problem}

\begin{solution}
\begin{enumerate}
\item 正方形 \(AA_{1}C_{1}C\) 给出 \(AA_{1}\perp AC\). 又 \(AC=4\)，故
\(AC^2+AB^2=4^2+3^2=25=BC^2\)，由勾股定理的逆定理得 \(AB\perp AC\). 两平面 \(ABC\) 与 \(AA_{1}C_{1}C\) 的交线为 \(AC\)，且两平面垂直；因此平面 \(ABC\) 内垂直于 \(AC\) 的直线 \(AB\) 垂直于平面 \(AA_{1}C_{1}C\)，从而 \(AB\perp AA_{1}\). 于是 \(AA_{1}\) 垂直于平面 \(ABC\) 内相交的两条直线 \(AB\) 和 \(AC\)，所以 \(AA_{1}\perp\) 平面 \(ABC\).

\item 建立空间直角坐标系，使 \(A\) 为原点，\(AB\)，\(AC\)，\(AA_{1}\) 分别为三个坐标轴的正方向，则
\(A=(0,0,0)\)，\(B=(3,0,0)\)，\(C=(0,4,0)\)，\(A_{1}=(0,0,4)\)，\(B_{1}=(3,0,4)\)，\(C_{1}=(0,4,4)\). 取
\(\overrightarrow{BC_{1}}=(-3,4,4)\)，\(\overrightarrow{BA_{1}}=(-3,0,4)\)，\(\overrightarrow{BB_{1}}=(0,0,4)\). 于是平面 \(A_{1}BC_{1}\) 和 \(B_{1}BC_{1}\) 的法向量可分别取为
\(\bs{n}_{1}=\overrightarrow{BC_{1}}\times\overrightarrow{BA_{1}}=(16,0,12)\) 和
\(\bs{n}_{2}=\overrightarrow{BC_{1}}\times\overrightarrow{BB_{1}}=(16,12,0)\). 因此
\(\bs{n}_{1}\cdot\bs{n}_{2}=256\)，且 \(\lvert\bs{n}_{1}\rvert=\lvert\bs{n}_{2}\rvert=20\). 设所求二面角为 \(\theta\)，则
\(\cos\theta=\dfrac{\lvert\bs{n}_{1}\cdot\bs{n}_{2}\rvert}{\lvert\bs{n}_{1}\rvert\lvert\bs{n}_{2}\rvert}=\dfrac{256}{20\cdot20}=\dfrac{16}{25}\).

\item 设 \(D=B+t\overrightarrow{BC_{1}}\)，其中 \(0\leqslant t\leqslant1\). 则
\(D=(3-3t,4t,4t)\)，所以 \(\overrightarrow{AD}=(3-3t,4t,4t)\)，而 \(\overrightarrow{A_{1}B}=(3,0,-4)\). 条件 \(AD\perp A_{1}B\) 等价于
\(\overrightarrow{AD}\cdot\overrightarrow{A_{1}B}=3(3-3t)-16t=9-25t=0\)，解得 \(t=\dfrac9{25}\). 因为 \(\dfrac9{25}\in[0,1]\)，所以在线段 \(BC_{1}\) 上存在这样的点 \(D\)，且
\(\dfrac{BD}{BC_{1}}=t=\dfrac9{25}\).
\end{enumerate}
\end{solution}

\begin{problem}
设 \(L\) 为曲线 \(C: y = \frac{\ln x}{x}\) 在点 \((1,0)\) 处的切线.
\begin{enumerate}
\item 求 \(L\) 的方程；
\item 证明：除切点 \((1,0)\) 之外，曲线 \(C\) 在直线 \(L\) 的下方.
\end{enumerate}
\end{problem}

\begin{solution}
\begin{enumerate}
\item 函数 \(f(x)=\dfrac{\ln x}{x}\) 的定义域为 \(x>0\)，且
\(f'(x)=\dfrac{1-\ln x}{x^2}\)，所以 \(f'(1)=1\). 切线经过 \((1,0)\)，故其方程为 \(L:y=x-1\).

\item 令 \(h(x)=x(x-1)-\ln x\)，其中 \(x>0\). 则
\(h'(x)=2x-1-\dfrac1x=\dfrac{(2x+1)(x-1)}x\). 当 \(0<x<1\) 时，\(h'(x)<0\)，所以 \(h\) 在 \((0,1)\) 上递减；当 \(x>1\) 时，\(h'(x)>0\)，所以 \(h\) 在 \((1,+\infty)\) 上递增. 又 \(h(1)=0\)，因此当 \(x>0\) 且 \(x\ne1\) 时，\(h(x)>0\).

因为 \(x>0\)，由 \(h(x)>0\) 可得 \(\dfrac{\ln x}{x}<x-1\)，即除切点外，曲线 \(C\) 在直线 \(L\) 的下方.
\end{enumerate}
\end{solution}

\begin{problem}
已知 \(A\)，\(B\)，\(C\) 是椭圆 \(W: \frac{x^2}{4} + y^2 = 1\) 上的三个点，\(O\) 是坐标原点.
\begin{enumerate}
\item 当点 \(B\) 是 \(W\) 的右顶点，且四边形 \(OABC\) 为菱形时，求此菱形的面积；
\item 当点 \(B\) 不是 \(W\) 的顶点时，判断四边形 \(OABC\) 是否可能为菱形，并说明理由.
\end{enumerate}
\end{problem}

\begin{solution}
\begin{enumerate}
\item 椭圆 \(W\) 的右顶点为 \(B=(2,0)\). 若四边形 \(OABC\) 为菱形，则 \(OB\) 和 \(AC\) 是两条互相垂直且互相平分的对角线. \(OB\) 的中点为 \((1,0)\)，且 \(OB\) 在 \(x\) 轴上，所以 \(AC\) 是直线 \(x=1\). 代入椭圆方程得
\(y^2=1-\dfrac14=\dfrac34\)，故可取 \(A=\left(1,\dfrac{\sqrt3}{2}\right)\)，\(C=\left(1,-\dfrac{\sqrt3}{2}\right)\). 于是 \(AC=\sqrt3\)，\(OB=2\)，菱形面积为
\(S=\dfrac12\cdot AC\cdot OB=\sqrt3\).

\item 反设 \(B=(u,v)\) 不是椭圆的顶点，但四边形 \(OABC\) 为菱形. 设两条对角线 \(OB\) 和 \(AC\) 的交点为 \(M\)，并设 \(\bs{d}=\overrightarrow{MA}=(d_x,d_y)\). 由对角线互相平分，
\(M=\left(\dfrac u2,\dfrac v2\right)\)，\(A=M+\bs{d}\)，\(C=M-\bs{d}\)；由对角线互相垂直，得
\(\bs{d}\cdot B=0\)，即 \(u d_x+v d_y=0\).

将 \(A\) 和 \(C\) 的坐标代入椭圆方程，分别得到
\(\dfrac{(\frac u2+d_x)^2}{4}+\left(\frac v2+d_y\right)^2=1\) 和
\(\dfrac{(\frac u2-d_x)^2}{4}+\left(\frac v2-d_y\right)^2=1\). 两式相减得
\(u d_x+4v d_y=0\). 将此式与 \(u d_x+v d_y=0\) 相减得 \(3v d_y=0\). 另一方面，将前面的两个椭圆方程相加并利用 \(B\) 在椭圆 \(W\) 上，即 \(\dfrac{u^2}{4}+v^2=1\)，可得
\(\dfrac{d_x^2}{4}+d_y^2=\dfrac34\)，所以 \(\bs{d}\ne\bs{0}\).

若 \(v\ne0\)，则 \(d_y=0\)，再由 \(u d_x+v d_y=0\) 及 \(d_x\ne0\) 得 \(u=0\)，于是 \(v=\pm1\)，这时 \(B\) 是椭圆的上、下顶点. 若 \(v=0\)，由 \(\dfrac{u^2}{4}=1\) 得 \(u=\pm2\)，这时 \(B\) 是椭圆的左、右顶点. 两种情形都与 \(B\) 不是顶点矛盾，因此四边形 \(OABC\) 不可能为菱形.
\end{enumerate}
\end{solution}

\begin{problem}
已知 \(\{a_{n}\}\) 是由非负整数组成的无穷数列，该数列前 \(n\) 项的最大值记为 \(A_{n}\)，第 \(n\) 项之后各项 \(a_{n+1}\)，\(a_{n+2}\)，\(\cdots\) 的最小值记为 \(B_{n}\)，\(d_{n} = A_{n} - B_{n}\).
\begin{enumerate}
\item 若 \( \{a_{n}\} \) 为 2, 1, 4, 3, 2, 1, 4, 3, \( \cdots \)，是一个周期为 4 的数列 （即对任意 \(n \in \N^{*}\)，\(a_{n+4} = a_{n}\) ），写出 \(d_{1}\)，\(d_{2}\)，\(d_{3}\)，\(d_{4}\) 的值；
\item 设 \(d\) 是非负整数，证明： \(d_{n} = -d \) （\(n = 1\)，\(2\)，\(3 \cdots\)）的充分必要条件为 \(\{a_{n}\}\) 是公差为 \(d\) 的等差数列；
\item 证明：若 \(a_{1} = 2\)，\(d_{n} = 1 \)（\(n = 1\)，\(2\)，\(3\)，\(\cdots\)），则 \(\{a_{n}\}\) 的项只能是 1 或者 2，且有无穷多项为 1.
\end{enumerate}
\end{problem}

\begin{solution}
\begin{enumerate}
\item 由周期性，后面各项的最小值均为 \(1\). 因此 \(A_1=2\)，\(A_2=2\)，\(A_3=4\)，\(A_4=4\)，而 \(B_1=B_2=B_3=B_4=1\). 所以
\(d_1=1\)，\(d_2=1\)，\(d_3=3\)，\(d_4=3\).

\item 充分性：若 \(\{a_n\}\) 是公差为 \(d\) 的等差数列，则 \(a_n=a_1+(n-1)d\). 由于 \(d\geqslant0\)，数列不减，故 \(A_n=a_n\)，且尾数列从 \(a_{n+1}\) 开始不减，所以 \(B_n=a_{n+1}\). 于是
\(d_n=A_n-B_n=a_n-a_{n+1}=-d\).

必要性：设对每个正整数 \(n\) 都有 \(d_n=-d\)，则 \(B_n=A_n+d\). 记
\(\Delta_n=A_{n+1}-A_n\geqslant0\)，于是 \(B_{n+1}=B_n+\Delta_n\). 根据 \(B_n\) 的定义，
\(B_n=\min\{a_{n+1},B_{n+1}\}\).

若 \(\Delta_n=0\)，则 \(a_{n+1}\leqslant A_n\)，且 \(B_{n+1}=B_n\). 由 \(B_n=\min\{a_{n+1},B_n\}\) 得 \(a_{n+1}\geqslant B_n=A_n+d\). 结合 \(a_{n+1}\leqslant A_n\) 和 \(d\geqslant0\)，只能有 \(d=0\) 且 \(a_{n+1}=A_n=A_n+d\).

若 \(\Delta_n>0\)，则 \(a_{n+1}>A_n\)，所以 \(a_{n+1}=A_n+\Delta_n\). 又
\(B_{n+1}=B_n+\Delta_n=A_n+d+\Delta_n=a_{n+1}+d\geqslant a_{n+1}\)，故 \(B_n=a_{n+1}\). 于是 \(a_{n+1}=B_n=A_n+d\)，并且 \(A_{n+1}=a_{n+1}\).

综上，对每个 \(n\) 均有 \(a_{n+1}=A_n+d\) 且 \(A_{n+1}=a_{n+1}\). 由 \(A_1=a_1\) 归纳得到 \(A_n=a_n\)，从而 \(a_{n+1}-a_n=d\). 所以 \(\{a_n\}\) 是公差为 \(d\) 的等差数列.

\item \(A_1=2\) 且 \(d_1=1\)，所以 \(B_1=1\)，从而 \(a_n\geqslant1\) 对所有 \(n\geqslant2\) 成立. 若存在 \(a_m\geqslant3\)，取第一个这样的下标 \(m\)，则 \(m\geqslant2\)，且 \(a_1\)，\(\ldots\)，\(a_{m-1}\) 均为 \(1\) 或 \(2\)，所以 \(A_{m-1}=2\)，\(B_{m-1}=1\). 因此在 \(a_m\) 之后存在某项等于 \(1\). 但 \(A_m\geqslant3\)，由 \(d_m=1\) 得 \(B_m=A_m-1\geqslant2\)，这意味着 \(a_{m+1}\)，\(a_{m+2}\)，\(\ldots\) 均不小于 \(2\)，与前述存在后续的 \(1\) 矛盾. 因此不存在 \(a_m\geqslant3\)，结合非负性可知每一项只能是 \(1\) 或 \(2\).

此时对任意 \(n\) 都有 \(A_n=2\)，由 \(d_n=1\) 得 \(B_n=1\). 若数列中只有有限多项为 \(1\)，取 \(n\) 大于最后一个 \(1\) 的下标，则 \(a_{n+1}\)，\(a_{n+2}\)，\(\ldots\) 全为 \(2\)，从而 \(B_n=2\)，矛盾. 因此数列中有无穷多项为 \(1\).
\end{enumerate}
\end{solution}
